\documentclass[../../main.tex]{subfiles}

\begin{document}

\section{Additional Fundamental Topics}

This section discusses frequently appearing concepts in matrices
that they have their own names!

\subsection{Transpose}

Let's begin with transpose of a matrix.

\begin{definition}{}{}
The \textbf{transpose} of a matrix $A$ is defined by a matrix that
has switched rows and columns of $A$, i.e. each element is defined
as following.
\[
(A^\intercal)_{ij} = A_{ji}
\]
\end{definition}

Let's take a look at an example. You could think of it as holding
the left top and bottom right of a matrix and rotating to see its
back. For matrix $A$ defined as the following,
\[
A =
\begin{bmatrix}
    1 & 2 & 3 & 4 \\
    5 & 6 & 7 & 8
\end{bmatrix}
\]
its transpose, or $A$ transpose is defined as the following.
\[
A^\intercal =
\begin{bmatrix}
    1 & 5 \\
    2 & 6 \\
    3 & 7 \\
    4 & 8
\end{bmatrix}
\]
As you can see, if $A$ has the order $m \times n$, its transpose
will have the order $n \times m$. Another observation that we can
make is that a transpose of a row matrix is a column matrix and
vice versa. Here is a quick example. Assume matrix $A$ defined
as the following.
\[
A =
\begin{bmatrix}
    1 & 2 & 3
\end{bmatrix}
\]
The transpose of $A$ can be represented as the following.
\[
A^\intercal =
\begin{bmatrix}
    1 \\
    2 \\
    3
\end{bmatrix}
\]
Notice that the transpose of transpose of $A$ is simply $A$.
\[
(A^\intercal)^\intercal =
\begin{bmatrix}
    1 & 2 & 3
\end{bmatrix}
= A
\]
Before formally stating our observations, here are few definitions
to note.

\begin{definition}{}{}
A positive integral powers of $A$, or multiplying the matrix $n$
times can be represented as $A^n$. $A^0$ is defined to be the
identity matrix $I_n$ with corresponding order $n$. A square
matrix is \textbf{nilpotent} if there exists $n \in \mathbb{Z}$
such that $A^n = 0$.
\end{definition}

Other definitions from transpose is symmetry and skew-symmetry.

\begin{definition}{}{}
For matrix $A$ and its transpose $A^\intercal$, matrix $A$ is
\textbf{symmetric} if $A^\intercal = A$ and \textbf{skew-symmetric}
if $A^\intercal = -A$.
\end{definition}

An example of a symmetric matrix is the identity matrix. For more
general example, matrix $A$ and $B$ below are symmetric and
skew-symmetric respectively.
\[
A =
\begin{bmatrix}
    a & b & c \\
    b & d & e \\
    c & e & f
\end{bmatrix},
\quad
B =
\begin{bmatrix}
    0  & a  & b \\
    -a & 0  & c \\
    -b & -c & 0
\end{bmatrix}
\]
Notice that all symmetric and skew-symmetric matrices are square
matrices. Moreover, the diagonal elements must be zero for a
square matrix to be skew-symmetric.

Now let's discuss the properties of transpose of a matrix.

\begin{theorem}{}{Properties of Transpose of a Matrix}
For matrix $A$, $B$, scalar $\lambda$ and positive integer $n$,
the following equations are true assuming that the addition and
multiplication conditions are satisfied.
\begin{enumerate}
\item $(A^\intercal)^\intercal = A$
\item $(A + B)^\intercal = A^\intercal + B^\intercal$
\item $(\lambda A)^\intercal = \lambda A^\intercal$
\item $(AB)^\intercal = B^\intercal A^\intercal$
\item $(A^n)^\intercal = (A^\intercal)^n$
\end{enumerate}
\end{theorem}

There are lots of properties to prove, but let's do them one by one.
Below is the proof for the first property.

\begin{proof}
Let $A = [a_{ij}]_{m \times n}$. By definition, its transpose
$A^\intercal \coloneqq [a_{ji}]_{n \times m}$. Continuing,
$\left( [a_{ji}]_{n \times m} \right)^\intercal =
[a_{ij}]_{m \times n} = A$. Therefore, $(A^\intercal)^\intercal = A$.
\end{proof}

Below is the proof for the second property.

\begin{proof}
Let $A = [a_{ij}]_{m \times n}$ and $B = [b_{ij}]_{m \times n}$.
Consider the following equation.
\begin{align*}
(A + B)^\intercal
&= \left(
    [a_{ij}]_{m \times n} + [b_{ij}]_{m \times n}
\right)^\intercal
= \left(
    [a_{ij} + b_{ij}]_{m \times n}
\right)^\intercal \\
&= [a_{ji} + b_{ji}]_{n \times m}
= [a_{ji}]_{n \times m} + [b_{ji}]_{n \times m} \\
&= A^\intercal + B^\intercal
\end{align*}
Thus, the property holds.
\end{proof}

Let's continue our proof for $(\lambda A)^\intercal =
\lambda A^\intercal$.

\begin{proof}
Let $A = [a_{ij}]_{m \times n}$. Because
$(\lambda A)^\intercal
= [\lambda a_{ij}]_{m \times n}^\intercal
= [\lambda a_{ji}]_{n \times m}
= \lambda [a_{ji}]_{n \times m}
= \lambda A^\intercal$, the property holds.
\end{proof}

We have two more left! Here is the proof for the fourth property.

\begin{proof}
Let $A = [a_{ij}]_{m \times n}$ and $B = [b_{ij}]_{n \times p}$.
The following equation is satisfied by definitions of matrix
multiplication and transpose.
\[
(AB)^\intercal
= \left(
    [a_{ij}]_{m \times n} [b_{ij}]_{n \times p}
\right)^\intercal
= \left(
    \left[ \sum_{k = 1}^{n} a_{ik} b_{kj} \right]_{m \times p}
\right)^\intercal
= \left[ \sum_{k = 1}^{n} a_{jk} b_{ki} \right]_{p \times m}
\]
Continuing from the right hand side, the following equation is
obtained.
\[
B^\intercal A^\intercal
= \left[ b_{ij}^\intercal \right]_{p \times n}
    \left[ a_{ij}^\intercal \right]_{n \times m}
= \left[
    \sum_{k = 1}^{n} b_{ik}^\intercal a_{kj}^\intercal
\right]_{p \times m}
= \left[ \sum_{k = 1}^{n} b_{ki} a_{jk} \right]_{p \times m}
= \left[ \sum_{k = 1}^{n} a_{jk} b_{ki} \right]_{p \times m}
\]
Because both the left-hand side and the right-hand side of the
equation lead to the same matrix, the property holds.
\end{proof}

Finally, below is the proof for the corollary of the fourth property.

\begin{proof}
First, notice that by the fourth property, the following equation
holds.
\[
\left( A^2 \right)^\intercal
= A^\intercal A^\intercal
= \left( A^\intercal \right)^2
\]
Therefore, it suffices to show that $(A^n)^\intercal =
(A^\intercal)^n$ holds for $n = k + 1$ if it is true for $n = k$.
Assume $(A^k)^\intercal = (A^\intercal)^k$. By the fourth property,
the following equation is obtained.
\[
\left( A^{k+1} \right)^\intercal
= (A^k A)^\intercal
= A^\intercal (A^k)^\intercal
= A^\intercal (A^\intercal)^k
= (A^\intercal)^{k + 1}
\]
By induction, the property holds for $k \geq 2$. For $k = 1$,
it is self evident that the property holds true. Therefore, the
property holds for all positive integers $n$.
\end{proof}

For the next part of the section, we will discuss partitions.

\subsection{Partitions}

As always, let's start with key definitions.

\begin{definition}{}{}
As the name suggests a \textbf{submatrix} is a matrix of $A$ that
is obtained by omitting certain rows or columns.
\end{definition}

Consider matrix $A$ defined as the following.
\[
A =
\begin{bmatrix}
    1 & 2 & 3 \\
    4 & 5 & 6 \\
    7 & 8 & 9
\end{bmatrix}
\]
Below are a few example of submatrices of $A$.
\begin{align*}
A =
\begin{bmatrix}
    \tikzmarknode{a11}{1} &
    \tikzmarknode{a12}{2} &
    \tikzmarknode{a13}{3} \\
    \tikzmarknode{a21}{4} &
    \tikzmarknode{a22}{5} &
    \tikzmarknode{a23}{6} \\
    \tikzmarknode{a31}{7} &
    \tikzmarknode{a32}{8} &
    \tikzmarknode{a33}{9}
\end{bmatrix}
\begin{tikzpicture}[
        overlay, remember picture,
        arrow/.style = {-, very thick, opacity=0.7}
    ]
    \draw[arrow, Red] (a21.west)  -- (a23.east);
\end{tikzpicture}
&\rightarrow
\begin{bmatrix}
    1 & 2 & 3 \\
    7 & 8 & 9
\end{bmatrix} \\
A =
\begin{bmatrix}
    \tikzmarknode{a11}{1} &
    \tikzmarknode{a12}{2} &
    \tikzmarknode{a13}{3} \\
    \tikzmarknode{a21}{4} &
    \tikzmarknode{a22}{5} &
    \tikzmarknode{a23}{6} \\
    \tikzmarknode{a31}{7} &
    \tikzmarknode{a32}{8} &
    \tikzmarknode{a33}{9}
\end{bmatrix}
\begin{tikzpicture}[
        overlay, remember picture,
        arrow/.style = {-, very thick, opacity=0.7}
    ]
    \draw[arrow, Red] (a11.west) -- (a13.east);
    \draw[arrow, Red] (a31.west) -- (a33.east);
    \draw[arrow, Red] (a12.north) -- (a32.south);
\end{tikzpicture}
&\rightarrow
\begin{bmatrix}
    4 & 6
\end{bmatrix}
\end{align*}

\begin{definition}{}{}
We say the matrix is \textbf{partitioned} if horizontal and vertical
lines are used to divide the matrix into submatrices.
\end{definition}

Here are a few examples of different partitions of $A$.
\[
\begin{bNiceArray}{ccc}[margin]
    1 & 2 & 3 \\
    4 & 5 & 6 \\
    \hline
    7 & 8 & 9
\end{bNiceArray},
\quad
\begin{bNiceArray}{c|cc}[margin]
    1 & \textcolor{Red}{2} & \textcolor{Red}{3} \\
    4 & \textcolor{Red}{5} & \textcolor{Red}{6} \\
    \hline
    7 & 8 & 9
\end{bNiceArray},
\quad
\begin{bNiceArray}{c|c|c}[margin]
    1 & 2 & 3 \\
    \hline
    4 & 5 & 6 \\
    7 & 8 & 9
\end{bNiceArray}
\]
From different partitions, we can obtain different \textbf{blocks}.

\begin{definition}{}{}
A \textbf{block} of a partition of $A$ is a submatrix of $A$ divided
by the horizontal and vertical partition lines.
\end{definition}

The highlighted portion of the second partition above is an example
of a block. Now, why should we even care about partitions and blocks?
Partitions of a matrix are important because it could be a very
useful tool in computationally heavy matrix multiplication.
Let's take a look at a generalized example $AB$.
\[
AB =
\begin{bmatrix}
    a_{11} & a_{12} & a_{13} & a_{14} \\
    a_{21} & a_{22} & a_{23} & a_{24} \\
    a_{31} & a_{32} & a_{33} & a_{34} \\
    a_{41} & a_{42} & a_{43} & a_{44}
\end{bmatrix}
\begin{bmatrix}
    b_{11} & b_{12} & b_{13} \\
    b_{21} & b_{22} & b_{23} \\
    b_{31} & b_{32} & b_{33} \\
    b_{41} & b_{42} & b_{43}
\end{bmatrix}
\]
At least for me, this multiplication looks very painful. Fortunately,
using partitions, we can turn this monster into a baby monster.
\begin{align*}
AB
&=
\begin{bNiceArray}{cc|cc}[margin]
    a_{11} & a_{12} & a_{13} & a_{14} \\
    a_{21} & a_{22} & a_{23} & a_{24} \\
    a_{31} & a_{32} & a_{33} & a_{34} \\
    a_{41} & a_{42} & a_{43} & a_{44}
\end{bNiceArray}
\begin{bNiceArray}{ccc}[margin]
    b_{11} & b_{12} & b_{13} \\
    b_{21} & b_{22} & b_{23} \\
    \hline
    b_{31} & b_{32} & b_{33} \\
    b_{41} & b_{42} & b_{43}
\end{bNiceArray} \\
&=
\begin{bmatrix}
    a_{11} & a_{12} \\
    a_{21} & a_{22} \\
    a_{31} & a_{32} \\
    a_{41} & a_{42}
\end{bmatrix}
\begin{bmatrix}
    b_{11} & b_{12} & b_{13} \\
    b_{21} & b_{22} & b_{23}
\end{bmatrix}
+
\begin{bmatrix}
    a_{13} & a_{14} \\
    a_{23} & a_{24} \\
    a_{33} & a_{34} \\
    a_{43} & a_{44}
\end{bmatrix}
\begin{bmatrix}
    b_{31} & b_{32} & b_{33} \\
    b_{41} & b_{42} & b_{43}
\end{bmatrix} \\
&=
\begin{bmatrix}
    a_{11}b_{11} + a_{12}b_{21} &
    a_{11}b_{12} + a_{12}b_{22} &
    a_{11}b_{13} + a_{12}b_{23} \\
    a_{21}b_{11} + a_{22}b_{21} &
    a_{21}b_{12} + a_{22}b_{22} &
    a_{21}b_{13} + a_{22}b_{23} \\
    a_{31}b_{11} + a_{32}b_{21} &
    a_{31}b_{12} + a_{32}b_{22} &
    a_{31}b_{13} + a_{32}b_{23} \\
    a_{41}b_{11} + a_{42}b_{21} &
    a_{41}b_{12} + a_{42}b_{22} &
    a_{41}b_{13} + a_{42}b_{23}
\end{bmatrix} \\
&\qquad\qquad\qquad\qquad+
\begin{bmatrix}
    a_{13}b_{31} + a_{14}b_{41} &
    a_{13}b_{32} + a_{14}b_{42} &
    a_{13}b_{33} + a_{14}b_{43} \\
    a_{23}b_{31} + a_{24}b_{41} &
    a_{23}b_{32} + a_{24}b_{42} &
    a_{23}b_{33} + a_{24}b_{43} \\
    a_{33}b_{31} + a_{34}b_{41} &
    a_{33}b_{32} + a_{34}b_{42} &
    a_{33}b_{33} + a_{34}b_{43} \\
    a_{43}b_{31} + a_{44}b_{41} &
    a_{43}b_{32} + a_{44}b_{42} &
    a_{43}b_{33} + a_{44}b_{43}
\end{bmatrix} \\
&=
\begin{bmatrix}
    \scriptstyle
    a_{11}b_{11} + a_{12}b_{21} + a_{13}b_{31} + a_{14}b_{41} &
    \scriptstyle
    a_{11}b_{12} + a_{12}b_{22} + a_{13}b_{32} + a_{14}b_{42} &
    \scriptstyle
    a_{11}b_{13} + a_{12}b_{23} + a_{13}b_{33} + a_{14}b_{43} \\
    \scriptstyle
    a_{21}b_{11} + a_{22}b_{21} + a_{23}b_{31} + a_{24}b_{41} &
    \scriptstyle
    a_{21}b_{12} + a_{22}b_{22} + a_{23}b_{32} + a_{24}b_{42} &
    \scriptstyle
    a_{21}b_{13} + a_{22}b_{23} + a_{23}b_{33} + a_{24}b_{43} \\
    \scriptstyle
    a_{31}b_{11} + a_{32}b_{21} + a_{33}b_{31} + a_{34}b_{41} &
    \scriptstyle
    a_{31}b_{12} + a_{32}b_{22} + a_{33}b_{32} + a_{34}b_{42} &
    \scriptstyle
    a_{31}b_{13} + a_{32}b_{23} + a_{33}b_{33} + a_{34}b_{43} \\
    \scriptstyle
    a_{41}b_{11} + a_{42}b_{21} + a_{43}b_{31} + a_{44}b_{41} &
    \scriptstyle
    a_{41}b_{12} + a_{42}b_{22} + a_{43}b_{32} + a_{44}b_{42} &
    \scriptstyle
    a_{41}b_{13} + a_{42}b_{23} + a_{43}b_{33} + a_{44}b_{43}
\end{bmatrix}
\end{align*}

Well, this looks good to me! If the matrices being multiplied
have higher dimension, partitioning such matrices could help in
computation and reducing mistakes. Now that we have seen partitions
of matrices, we could derive special forms from the ideas.

\subsection{Special Forms of Matrix}

As always, let's build from definitions.

\begin{definition}{}{}
A \textbf{symmetrically partitioned matrix} refers to a partitioned
matrix such that its arrangements of the blocks are diagonally
symmetric. The \textbf{diagonal blocks} are the blocks that pass
through the diagonal symmetry.
\end{definition}

We could also define matrices with diagonal blocks that are not
necessarily symmetrically partitioned.

\begin{definition}{}{}
A \textbf{block diagonal matrix} is a matrix with square diagonal
blocks and zero on the other entries.
\end{definition}

General form for square matrices $A_i$ can be represented as $A =
\begin{bmatrix}
    A_1    & 0      & \cdots & 0 \\
    0      & A_2    & \cdots & 0 \\
    \vdots & \vdots & \ddots & 0 \\
    0      & 0      & 0      & 0 \\
\end{bmatrix}$.
Continuing, we could define rows and columns based on its elements.

\begin{definition}{}{}
\textbf{Zero rows} and \textbf{zero columns} are rows and columns
composed of zeros respectively.
\end{definition}

This definition is quite straightforward, and now that we have
basic terms in mind, let's discuss two forms from the definitions
above.

\begin{definition}{}{}
Matrix in \textbf{row echelon form (REF)} satisfies the conditions
below.
\begin{itemize}[noitemsep]
\item
All nonzero leading entry is $1$. (Side note here is that some
textbooks and lectures do not require the leading entry to be $1$,
but I actually would like to keep it this way as this is how I
initially learned it.)
\item
All nonzero leading entries in a nonzero row in a later row are to
the right of nonzero leading entries in previous rows.
\item
If there are both nonzero and zero rows, all zero rows must appear
after the nonzero rows above.
\end{itemize}

The leading entry of matrices in REF is known as a \textbf{pivot}.
Similarly, matrices in \textbf{reduced row echelon form (RREF)}
are matrices that satisfy the following conditions.
\begin{itemize}[noitemsep]
\item
All nonzero leading entries in a nonzero row is $1$.
\item
If there are both nonzero and zero rows, all zero rows must appear
after the nonzero rows above.
\item
All leading $1$s in the succeeding rows appear in a column after the
columns with previous leading $1$s.
\item
The leading $1$ must be the only nonzero element present in each
column if there exists such $1$.
\end{itemize}
\end{definition}

Below are examples of a matrices in REF and RREF respectively.
\[
\begin{bmatrix}
    1 & 2 & 3 & 4 \\
    0 & 0 & 1 & 2 \\
    0 & 0 & 0 & 1 \\
    0 & 0 & 0 & 0
\end{bmatrix},
\quad
\begin{bmatrix}
    1 & 0 & 2 & 0 \\
    0 & 1 & 2 & 0 \\
    0 & 0 & 0 & 1 \\
    0 & 0 & 0 & 0
\end{bmatrix}
\]
Matrices that follow these forms will be very useful when solving
linear systems of equations, but we will save that for the latter
notes. Now let's discuss about triangular matrices.

Triangular matrices are a special type of square matrices that is
determined by elements below the main diagonal. If all its
elements below the main diagonal are zeros, we call it an
\textbf{upper triangular} matrix. Similarly, if all elements above
the main diagonal are zeros, we call them \textbf{lower triangular}
matrices. For more formal statement, consider the following
definition.

\begin{definition}{}{}
A square matrix $A = [a_{ij}]$ is \textbf{upper triangular} if
$a_{ij} = 0 \, \forall \, i > j$. Similarly, the matrix is
\textbf{lower triangular} if $a_{ij} = 0 \, \forall \, i < j$.
\end{definition}

Let's take a look at an example. The left and right matrices below
are upper triangular and lower triangular matrices respectively.
\[
\begin{bmatrix}
    1 & 2 & 3 & 4 \\
    0 & 2 & 3 & 4 \\
    0 & 0 & 3 & 4 \\
    0 & 0 & 0 & 4
\end{bmatrix},
\quad
\begin{bmatrix}
    1 & 0 & 0 & 0 \\
    2 & 1 & 0 & 0 \\
    3 & 2 & 1 & 0 \\
    0 & 3 & 2 & 0
\end{bmatrix}
\]
Now that we have definitions of triangular matrices in mind,
we can discuss a theorem from such matrices.

\begin{theorem}{}{Property of Triangular Matrix Multiplication}
The product of two lower triangular matrices is a lower triangular
matrix, and the product of two upper triangular matrices is an upper
triangular matrix.
\end{theorem}
\begin{proof}
Let $A = [a_{ij}]_{n \times n}$ and $B = [b_{ij}]_{n \times n}$ be
lower triangular matrices. Consider the following equation for their
product matrix $C = [c_{ij}]_{n \times n}$.
\[
C
= \left[ \sum_{k = 1}^{n} a_{ik} b_{kj} \right]
= \left[
    \sum_{k = 1}^{j-1} a_{ik} b_{kj}
        + \sum_{k = j}^{n} a_{ik} b_{kj}
\right]
\]
From the expression, it suffices to show that $\sum_{k = 1}^{j-1}
a_{ik} b_{kj} + \sum_{k = j}^{n} a_{ik} b_{kj} = 0$ for $i < j$.
Notice that by definition, $b_{kj} = 0$ for all $k = 1$ to $k = j-1$.
Similarly, $a_{ik} = 0$ for all integer $k \in [j, n]$. Therefore,
$c_{ij} = 0$ for all $i < j$, proving the first half of the theorem.

Similarly, let $A' = [a'_{ij}]_{n \times n}$ and $B' =
[b'_{ij}]_{n \times n}$ be upper triangular matrices. The product
matrix $C'= [c'_{ij}]_{n \times n}$ can be represented as following.
\[
C'
= \left[ \sum_{k = 1}^{n} a'_{ik} b'_{kj} \right]
= \left[
    \sum_{k = 1}^{j-1} a'_{ik} b'_{kj}
        + \sum_{k = j}^{n} a'_{ik} b'_{kj}
\right]
\]
By definition, it is evident that $a'_{ik} = 0$ for integer $k
\in [1, j-1]$ and $b'_{kj} = 0$ for integer $k \in [j, n]$.
Therefore, for $i > j$, $c'_{ij} = 0$ and the theorem is true.
\end{proof}

This is it for different topics on fundamentals of matrices! I know
this note is quite long, but let's conclude our notes with the
discussion on determinants.

\end{document}


